% section: 指数対数

  \section{指数対数}\label{節：指数対数}
  \noindent\textbf{【本編で使う場面】}\par
  指数法則と,積・商・累乗を和・差・定数倍に変える\bookterm{logarithm}{対数}法則を確認する.\sectionref*{節：対数型}では,真数の正値と底の条件も使う.

    \subsection{指数}\label{節：指数}
      $a$が$p$回掛けられている状況を考える.
      ここで
      \[
        a\times a\times\cdots\times a
      \]
      のように記すのは煩雑なので,
      \[
        a\times a\times\cdots\times a=a^p
      \]
      と表す.
      このときの$p$を指数といい,
      \[
        f(x)=a^x
      \]
      で表されるような関数を指数関数という.
      ここでは$p$を正の整数として導入した.底$a>0$の場合は,\,$a^0=1,\ a^{-p}=\dfrac1{a^p}$と定め,さらに有理数・実数の指数へと拡張できる.
      以降,実数の指数を用いるときは底$a>0$とする.
    \subsection{対数}\label{節：対数}
      $a^p=q$のとき,\,$p$は
      \[
        \log_a{q}=p
      \]
      のように表される.
      この時$\log$をログと読み,\,$p$を\bookterm{logarithm}{対数}という.
      また,\,$a$を底,\,$q$を真数と言う.
      一般に底$a$は$a>0,\,a\ne1$を満たす必要がある.
      また,
      \[
        e=\lim_{x \to \infty} \left(1+\frac{1}{x}\right)^x
      \]
      で定義される数$e$を自然\bookterm{logarithm}{対数}の底またはネイピア数という.
      ただし,以下の公式では底$a,c$は正で$1$ではなく,真数$p,M,N,b$は正とする.
      \bookterm{logarithm}{対数}に関して,一般に以下が成り立つ.
      \begin{align*}
        &\log_a{p^q}=q\log_a{p}\\
        &\log_a{MN}=\log_a{M}+\log_a{N}\\
        &\log_a{\frac{M}{N}}=\log_a{M}-\log_a{N}\\
        &\log_a{b}=\frac{\log_c{b}}{\log_c{a}}
      \end{align*}

      また,
      \[
        f(x)=\log_a x
      \]
      というような関数を\bookterm{logarithm}{対数}関数といい,指数関数の\bookterm{inverse}{逆関数}である.
      \bookterm{inverse}{逆関数}とは
      \[
        y=f(x)
      \]
      を満たすような$x,y$について,
      \[
        x=g(y)
      \]
      となるような関数$g$を指す.
      グラフ上では,\bookterm{inverse}{逆関数}は元の関数と関数$y=x$に対して対称である.
      関数$g$はよく,\,$f^{-1}$と書かれる.
      \bookterm{inverse}{逆関数}はある$x$と$y$が関数$f$で結ばれているという対応付けをそのままに,結び方を$y$について整理したものから$x$について整理したものに変えた,ということである.
